
\begin{center}
	\LARGE
	Uden hjælpemidler
\end{center}
\begin{opgavetekst}{Opgave 1}
	En differentialligning er givet ved
	\begin{align*}
		\frac{dy}{dt} = \frac{3t^2}{y}.
	\end{align*}
\end{opgavetekst}
\begin{delopgave}{1}
	Bestem $y'$ i punktet $P(2,3)$.
\end{delopgave}
\begin{delopgave}{2}
	Vis, at funktionen
	\begin{align*}
		f(t) = \sqrt{2t^3 + 4}
	\end{align*}
	er en partikulær løsning til differentialligningen. 
\end{delopgave}
%
%
%
%
%
%
\begin{opgavetekst}{Opgave 2}
	En differentialligning er givet ved
	\begin{align*}
		y' = x - 5y^2.
	\end{align*}
\end{opgavetekst}
\begin{delopgave}{1}
	Bestem linjeelementet i $P(-2,1)$.	
\end{delopgave}
\begin{delopgave}{2}
	Bestem ligningen for tangenten til løsningen af differentialligningen, der går gennem punktet $P(-2,1)$.
\end{delopgave}
%
%
%
%
%
\begin{opgavetekst}{Opgave 3}
	En differentialligning er givet ved
	\begin{align*}
		y' = 10 - 5y.
	\end{align*}
\end{opgavetekst}
\begin{delopgave}{1}
	Bestem den fuldstændige løsning til differentialligningen.
\end{delopgave}
\begin{delopgave}{2}
	Bestem den partikulære løsning til differentialligningen, der går gennem punktet $P(0,6)$.
\end{delopgave}
%
%
%
%
%
\newpage
\begin{opgavetekst}{Opgave 4}
	En differentialligning er givet ved
	\begin{align*}
		\frac{dy}{dx} = 2y(10-y).
	\end{align*}
\end{opgavetekst}
\begin{delopgave}{1}
	Bestem en partikulær løsning til differentialligningen, der går gennem punktet $P(0, 3)$.
\end{delopgave}
%
%
%
%
%
\begin{opgavetekst}{Opgave 5}
	Hældningsfeltet for en differentialligning er givet på Figur \ref{fig:opgfelt}.
	\begin{figure}[H]
	\centering
	\begin{tikzpicture}
	\begin{axis}[
		xmin = -5, xmax = 5,
		ymin = -0, ymax = 10,
		axis lines = center,
		xlabel = $x$,
		ylabel = $y$,
	]
		\foreach \xVal in {-5,-4.5,...,5}{
			\foreach \yVal in {-0,0.5,...,10}{
				\pgfmathsetmacro{\sizescale}{0.15};
				\pgfmathsetmacro{\slope}{-0.5*\yVal};
				\pgfmathsetmacro{\angle}{atan(\slope)};
				\pgfmathsetmacro{\xlen}{\sizescale*cos(\angle)};
				\pgfmathsetmacro{\ylen}{\sizescale*sin(\angle)};
				\edef\temp{\noexpand\draw [color = SeaGreen] 
				(axis cs: \xVal-\xlen, \yVal-\ylen) -- (axis cs: \xVal + \xlen, \yVal + \ylen);}
				\temp
			}
		}
	\end{axis}
	\end{tikzpicture}
	\caption{Hældningsfelt.}
	\label{fig:opgfelt}
	\end{figure}
	Det oplyses, at differentialligningen tilsvarer én af tre følgende differentialligninger.
	\begin{align*}
		l_1:& \ y' = 2y,\\
		l_2:& \ y' = -0.5y,\\
		l_3:& \ y' = 0.5y(10-y).
	\end{align*}
\end{opgavetekst}
\begin{delopgave}{1}
	Afgør, hvilken af de tre differentialligninger hældningsfeltet tilsvarer. Begrund dit svar. 
\end{delopgave}
%
%
%
%
%
\newpage
\begin{opgavetekst}{Opgave 6}
	En differentialligning er repræsenteret ved grafen på Figur \ref{fig:yymaerke}.
	\begin{figure}[H]
		\centering
		\begin{tikzpicture}
		\begin{axis}[
			axis lines = center, 
			xmin = -1.5, xmax = 8.5, 
			ymin = -1.5, ymax = 8.5,
			xlabel = $y$, ylabel = $y'$,
			grid = both,
			xtick = {-2,0,2,...,8},
			ytick = {-2,0,2,...,8},
			minor tick num = 1,
		]
			\addplot[color = SeaGreen, domain = -2:9, thick]
			 {-2*x + 8};
		\end{axis}
		\end{tikzpicture}
		\caption{Grafen for differentialligningen.}
		\label{fig:yymaerke}
	\end{figure}
	\phantom{h}
\end{opgavetekst}
\begin{delopgave}{1}
	Bestem den partikulære løsning til differentialligningen, der er repræsentet ved grafen på Figur \ref{fig:yymaerke}, og som går gennem punktet $P(0,5)$.
\end{delopgave}
%
%
%
%
%

\begin{opgavetekst}{Opgave 7}
	En differentialligning er repræsenteret ved grafen på Figur \ref{fig:yymaerke2}.
	\begin{figure}[H]
		\centering
		\begin{tikzpicture}
		\begin{axis}[
			axis lines = center, 
			xmin = -1.5, xmax = 8.5, 
			ymin = -1.5, ymax = 8.5,
			xlabel = $y$, ylabel = $y'$,
			grid = both,
			xtick = {-2,0,2,...,8},
			ytick = {-2,0,2,...,8},
			minor tick num = 1,
		]
			\addplot[color = SeaGreen, domain = -2:9, thick, samples = 200]
			 {0.5*x*(8-x)};
		\end{axis}
		\end{tikzpicture}
		\caption{Grafen for differentialligningen.}
		\label{fig:yymaerke2}
	\end{figure}
	\phantom{h}
\end{opgavetekst}
\begin{delopgave}{1}
	Bestem den fuldstændige løsning til differentialligningen, der er repræsenteret ved grafen på Figur \ref{fig:yymaerke2}.
\end{delopgave}
%
%%%%%%%%%%%%%%%
%
%
%
\begin{opgavetekst}{Opgave 8}
	Grafen for en harmonisk svingning $f$ givet ved
	\begin{align*}
		f(t) = A \sin(\omega t + \varphi) + B
	\end{align*}
	kan ses på Figur \ref{fig:opgharmonisk}.
	\begin{figure}[H]
		\centering
		\begin{tikzpicture}
		\begin{axis}
			[
			axis lines = center,
			xmin = -10, xmax = 10,
			ymin = -10, ymax = 6, 
			ytick = {-12,-10,...,10,12},
			xtick = {0},		
			]
			\addplot[color = teal, samples = 1000, domain = -20:20, thick] {5*sin(deg(x+ 3*3.1415/2)) - 1};
			\draw (axis cs: {3.1415 * (-6)} , 0.1) -- (axis cs: {3.1415 * (-6)} , -0.1)
			node[pos = 1, anchor = north] {$-6\pi$};
		%	\draw (axis cs: {3.1415 * (-5)} , 0.1) -- (axis cs: {3.1415 * (-5)} , -0.1)
		%	node[pos = 1, anchor = north] {$-5\pi$};			
			\draw (axis cs: {3.1415 * (-4)} , 0.1) -- (axis cs: {3.1415 * (-4)} , -0.1)
			node[pos = 1, anchor = north] {$-4\pi$};
		%	\draw (axis cs: {3.1415 * (-3)} , 0.1) -- (axis cs: {3.1415 * (-3)} , -0.1) 
		%	node[pos = 1, anchor = north] {$-3\pi$};
			\draw (axis cs: {3.1415 * (-2)} , 0.1) -- (axis cs: {3.1415 * (-2)} , -0.1)
			node[pos = 1, anchor = north] {$-2\pi$};
		%	\draw (axis cs: {3.1415 * (-1)} , 0.1) -- (axis cs: {3.1415 * (-1)} , -0.1)
		%	node[pos = 1, anchor = north] {$-\pi$};
			\draw (axis cs: {3.1415 * (1)} , 0.1) -- (axis cs: {3.1415 * (1)} , -0.1)
			node[pos = 1, anchor = north] {$\pi$};
			\draw (axis cs: {3.1415 * (2)} , 0.1) -- (axis cs: {3.1415 * (2)} , -0.1)
			node[pos = 1, anchor = north] {$2\pi$};
			\draw (axis cs: {3.1415 * (3)} , 0.1) -- (axis cs: {3.1415 * (3)} , -0.1)
			node[pos = 1, anchor = north] {$3\pi$};
			\draw (axis cs: {3.1415 * (4)} , 0.1) -- (axis cs: {3.1415 * (4)} , -0.1)
			node[pos = 1, anchor = north] {$4\pi$};
			\draw (axis cs: {3.1415 * (5)} , 0.1) -- (axis cs: {3.1415 * (5)} , -0.1)
			node[pos = 1, anchor = north] {$5\pi$};
			\draw (axis cs: {3.1415 * (6)} , 0.1) -- (axis cs: {3.1415 * (6)} , -0.1)
			node[pos = 1, anchor = north] {$6\pi$};
		\end{axis}
		\end{tikzpicture}
	\end{figure}
	\phantom{h}
\end{opgavetekst}
\begin{delopgave}{1}
	Bestem værdierne for $A$, $\omega$, $\varphi$ og $B$ samt perioden $T$ for den harmoniske svingning $f$.
\end{delopgave}	
%
%
%
%
%
\begin{opgavetekst}{Opgave 9}
\end{opgavetekst}
\begin{delopgave}{1}
	Løs ligningen 
	\begin{align*}
		\sin(x)^2 - \frac{\sqrt{2}}{2}\sin(x) = 0
	\end{align*}
	for $x \in [0, 2\pi]$. 
\end{delopgave}
%
%
%
%
%
%
\begin{opgavetekst}{Opgave 10}
\end{opgavetekst}
\begin{delopgave}{1}
	Løs integralet
	\begin{align*}
		\int_0^{\sqrt{\frac{\pi}{2}}} 6x\cos(x^2) dx.
	\end{align*}
\end{delopgave}
